\section{Problem Statement}\label{sec:problem}

This paper aims to leverage autonomous agents to automatically construct predictive algorithms on intensive care unit (ICU) time-series data that achieve both high predictive accuracy and clinical plausibility. Toward this goal, we first introduce the clinical prediction task specification, the candidate solution space, and the clinical feasibility constraint, followed by the formal problem definition and three key optimization challenges.

\textbf{Clinical Prediction Task Specification.} Consider a clinical time-series cohort $\mathcal{D}$ comprising $N$ ICU patient stays. Let $\mathcal{C}$ denote the set of clinical feature types, consisting of $d_s$ static demographic features and $d_d$ dynamic physiological and laboratory features with $|\mathcal{C}| = d_s + d_d$. For each patient stay $i \in \{1, \dots, N\}$, the observed record comprises a static demographic feature vector $\mathbf{x}_i^{\text{stat}} \in \mathbb{R}^{d_s}$ and an irregularly sampled multivariate dynamic time series $\mathcal{X}_i^{\text{dyn}} = \{(\tau,\, \mathbf{x}_{i,\tau}^{\text{dyn}})\}_{\tau \in \mathcal{T}_i^{\text{obs}}}$, where $\tau \ge 0$ is the timestamp since admission, $\mathbf{x}_{i,\tau}^{\text{dyn}} \in (\mathbb{R} \cup \{\bot\})^{d_d}$ is the observation vector at timestamp $\tau$, and $\bot$ denotes unmeasured missing entries. Let $\mathcal{X}_{i, \le t}^{\text{dyn}} = \{(\tau,\, \mathbf{x}_{i,\tau}^{\text{dyn}}) \in \mathcal{X}_i^{\text{dyn}} \mid 0 \le \tau \le t\}$ denote the observed dynamic history of patient $i$ up to time $t$. Depending on the prediction trigger mechanism and supervision structure, we study two classes of clinical prediction tasks with distinct temporal modeling complexities:
\begin{itemize}
\item \textbf{Stay-Level One-Shot Prediction:} Each stay is associated with a single outcome label $y_i$. A nonrolling prediction is triggered once at the end of a fixed observation window $W$ after admission, mapping the observed history within this fixed window to a scalar outcome prediction:
\begin{equation*}
f_{\text{stay}} : \big(\mathbf{x}_i^{\text{stat}},\, \mathcal{X}_{i, \le W}^{\text{dyn}}\big) \longmapsto \hat{y}_i.
\end{equation*}
\item \textbf{Hourly Rolling Prediction:} Each stay lasting $T_i$ hours is associated with a variable-length target sequence $(y_{i,1}, y_{i,2}, \dots, y_{i, T_i})$. As the decision time $t$ advances hourly, the model outputs a rolling prediction sequence $(\hat{y}_{i,1}, \hat{y}_{i,2}, \dots, \hat{y}_{i, T_i})$ for early warning of acute clinical events or real-time monitoring of patient status. To prevent temporal leakage caused by accessing future records with $\tau > t$, feature construction and prediction at any time $t$ depend strictly on the observed history $\mathcal{X}_{i, \le t}^{\text{dyn}}$ up to the current time:
\begin{equation*}
\hat{y}_{i,t} = f_{\text{roll}}\big(\mathbf{x}_i^{\text{stat}},\, \mathcal{X}_{i, \le t}^{\text{dyn}}\big), \quad t = 1, 2, \dots, T_i.
\end{equation*}
\end{itemize}

In terms of output formulation, both temporal task classes encompass classification and regression: classification tasks predict an event risk probability $\hat{y} \in [0, 1]$ for a binary target label $y \in \{0, 1\}$, whereas regression tasks predict a continuous physiological value $\hat{y} \in \mathbb{R}$ for a real-valued target $y \in \mathbb{R}$.

\textbf{Candidate Solution Space.} Unlike traditional modeling approaches that rely on manual trial-and-error by human experts, a large language model (LLM) agent operates within an isolated code execution environment and autonomously evolves predictive algorithms through multi-round closed-loop interactions of hypothesis generation, code synthesis, model training, and feedback observation. In each round, the agent generates or modifies an executable candidate solution $s = (g, \phi) \in \mathcal{G} \times \Phi$ to instantiate the prediction function $f \in \{f_{\text{stay}}, f_{\text{roll}}\}$. Specifically, the feature engineering module $\phi \in \Phi$ performs missing-value imputation, feature normalization, and temporal feature construction, while the predictive modeling module $g \in \mathcal{G}$ specifies the model family, network architecture, loss function, and training configuration to map the extracted features to predictions $\hat{y}_{i,t} = g\big(\phi(\mathbf{x}_i^{\text{stat}},\, \mathcal{X}_{i, \le t}^{\text{dyn}})\big)$, where stay-level one-shot prediction sets $t = W$. Given patient-disjoint training and validation splits $\mathcal{D}_{\text{train}}$ and $\mathcal{D}_{\text{val}}$, the execution environment trains $s$ on $\mathcal{D}_{\text{train}}$ and evaluates its validation performance $P_{\text{val}}(g, \phi) \in \mathbb{R}$ on $\mathcal{D}_{\text{val}}$.

\textbf{Clinical Feasibility Constraint.} On finite-sample clinical observational data, unconstrained search guided solely by maximizing a single scalar metric $P_{\text{val}}(g, \phi)$ is prone to exploiting single-feature autoregressive shortcuts or nonphysiological sampling artifacts that inflate validation scores. Such shortcuts disregard the pathophysiological mechanisms of multi-organ co-evolution and degrade model generalization under distribution shifts. To ensure pathophysiological plausibility, for a given clinical prediction task, let $\mathcal{C}^* \subset \mathcal{C}$ denote the subset of core features established by clinical guidelines, $\mathcal{M}_{\text{clin}}(g, \phi; \mathcal{C}^*)$ denote a consistency measure between the candidate solution and clinical guideline mechanisms, and $\Omega_{\text{clin}}$ denote the clinically acceptable feasible region. This feasible region requires the model to substantively rely on the guideline core features $\mathcal{C}^*$, avoid single-feature shortcut monopoly, and balance sensitivity and specificity in classification risk screening.

\textbf{Problem 1 (Clinically Constrained Agentic Clinical Algorithm Discovery):} \textit{Given a training set $\mathcal{D}_{\text{train}}$, a validation set $\mathcal{D}_{\text{val}}$, a clinical guideline core feature subset $\mathcal{C}^* \subset \mathcal{C}$, and a maximum candidate evaluation budget $B \in \mathbb{N}^+$, let $\mathcal{T} \subset \mathcal{G} \times \Phi$ denote the set of candidate solutions generated and evaluated by the agent over multi-round closed-loop interactions. The goal is to guide the evolution of $\mathcal{T}$ to identify an optimal candidate solution $s^* = (g^*, \phi^*)$ satisfying both the evaluation budget and the clinical feasibility constraint:}
\begin{equation}\label{eq:prob_def}
\begin{aligned}
s^* = (g^*, \phi^*) &= \operatorname*{arg\,max}_{(g, \phi) \in \mathcal{T}} \; P_{\text{val}}(g, \phi) \\
\text{s.t.} \quad & \mathcal{M}_{\text{clin}}(g, \phi; \mathcal{C}^*) \in \Omega_{\text{clin}}, \quad |\mathcal{T}| \le B.
\end{aligned}
\end{equation}

Solving \eqref{eq:prob_def} poses three key challenges:
\begin{itemize}
\item \textbf{Vast Code Search Space and Coarse-Grained Evaluation Feedback (C1):} Candidate solutions consist of free-form program code, yielding a vast search space. Moreover, a single proposed revision often modifies multiple coupled components simultaneously, whereas each validation run returns only an aggregate scalar score $P_{\text{val}}(g, \phi)$ for the entire solution, making it difficult to isolate the individual contribution of each modification.
\item \textbf{Guideline Alignment under Sparse Missingness and Single-Feature Shortcuts (C2):} On the one hand, real-world ICU records are irregularly sampled, and certain clinically informative features recommended by guidelines exhibit severe missingness in practice; rigidly enforcing guidelines by retaining high-missingness features regardless of empirical data availability introduces substantial noise. On the other hand, even when the core guideline features $\mathcal{C}^*$ are included in code, a black-box model may still rely on a single feature as a predictive shortcut, making code inspection alone insufficient to verify clinical plausibility.
\item \textbf{Search Interference and Computational Overhead under Coupled Dual-Space Alternation (C3):} The effectiveness of a feature transformation $\phi \in \Phi$ depends strongly on the inductive bias of the downstream model architecture $g \in \mathcal{G}$, preventing the problem from being decomposed into two independent subproblems. However, frequently alternating search between model architectures and feature transformations causes mutual interference and optimization instability, while increasing both search and computational costs.
\end{itemize}
